% ==========================================
% CHAPTER 5
% ==========================================
\chapter{Data I/O (Input/Output)} % 

\section{Loading and Writing Files} % 
At numerous points in this text we have alluded to the ability for loops to aid in the process of loading multiples files of data. %  
Now that you know how to concatenate strings and generate for-loops, we can cover the file loading/writing process. %  
There are several ways of opening data files in python. %  
Python itself has a built in mechanism for opening/writing files, and numpy also has support for file handling. %  
To open a file in python's interface, we type: % 

\begin{lstlisting}
>> file1 = open('filename.txt','w')
\end{lstlisting} % 

where \texttt{'w'} indicates we plan to write to the file. %  
(We could instead use \texttt{'r'} for read only, or \texttt{'a'} for appending to a file that already contains data). % 

\textbf{Example 5.1 Writing to a File} % 
\begin{lstlisting}
>> file1 = open('file.txt','w')
>> file1.write('this is a file')
>> file1.write('this is not a drill')
>> file1.close()
\end{lstlisting} % 

The close statement above tells python to close and save the file to the hard disk. % 

\textbf{Exercise 5.1 Writing data to a file} % 
\begin{lstlisting}
for i in range(len(planet_dist)-1):
    file.write(planet_dist[i] + planet_vel[i] + '\n')
\end{lstlisting} % 

where the \texttt{\textbackslash n} is necessary for us to create a new line when writing the file so the data will be properly divided into their respective row and column. % 

Numpy also has a file input output framework that is often useful to use. %  
The two we will discuss here are \texttt{np.loadtxt} and \texttt{np.genfromtxt}. %  

\textbf{Example 5.2 Loading files using loadtxt} % 
\begin{lstlisting}
data = np.loadtxt('filename.txt')
data = np.transpose(data)
times = data[0]
positions = data[1]
velocities = data[2]
\end{lstlisting} % 

Oftentimes data files have headers and footers- text that tells you what data is stored in the file. %  
For example, to skip a 5 line header and 3 line footer text, use: % 

\begin{lstlisting}
[IN]: data = np.genfromtxt('file.txt', skip_header=5, skip_footer=3)
\end{lstlisting} % 

\section{Loading Astronomical Fits Files} % 
In astronomy, in particular, it is useful to work with FITS (flexible image transport system) files. %  
FITS files are widely prevalent because of a feature they contain called a header. %  
Headers often contain information about the image itself. % 

To use them, we will have to import a library called pyfits (or on the lab computers, \texttt{astropy.io.fits}); %  
The syntax for opening a fits file is: % 

\begin{lstlisting}
>> hdu = pf.open(path)
\end{lstlisting} % 

\subsection{The Header} % 
The header is a dictionary containing a lot of useful information about the images stored in the fits file. %  

\begin{lstlisting}
>> head = hdu[0].header
\end{lstlisting} % 

\textbf{Example 5.3 Pulling from a header} % 
\begin{lstlisting}
ra = head['RA']
dec = head['DEC']
time = head['EXPTIME']
\end{lstlisting} % 

\subsection{The Image} % 
To access the image itself, we call: % 

\begin{lstlisting}
>> img = hdu[0].data
\end{lstlisting} % 
